% ---------------------------------------------------------------
% metadados.tex
% Editar APENAS este ficheiro para personalizar a capa do relatório.
% ---------------------------------------------------------------

\newcommand{\instituicao}{Instituto Superior de Engenharia do Porto}
\newcommand{\instituicaosigla}{ISEP}
\newcommand{\curso}{Mestrado em Engenharia de Inteligência Artificial}
\newcommand{\cursosigla}{MEIA}
\newcommand{\uc}{Engenharia do Conhecimento}
\newcommand{\ucsigla}{ENGCIA}
\newcommand{\anoletivo}{2026/2027}
\newcommand{\semestre}{1.\textordmasculine{} Semestre}

\newcommand{\tituloprincipal}{Sistema Pericial de Apoio à Decisão em Veículos Autónomos}
\newcommand{\subtitulo}{Relatório Prévio de Aquisição de Conhecimento}
\newcommand{\descricao}{Aquisição e representação de conhecimento pericial\\
para um sistema baseado em regras em Drools}

% ---------------------------------------------------------------
% Autoria
%
% Para acrescentar, remover ou reordenar autores, basta editar a lista
% abaixo: uma linha \autor{nome}{email} por elemento. A capa
% acompanha automaticamente, seja qual for o número de autores.
% ---------------------------------------------------------------
\providecommand{\autor}[2]{}

\newcommand{\listaautores}{%
  \autor{Ana Salomé Ferreira Teixeira}{1210803@isep.ipp.pt}%
  \autor{Diogo Miguel Gonçalves Ribeiro}{1260429@isep.ipp.pt}%
  \autor{João Gabriel Domingos de Azevedo}{1211519@isep.ipp.pt}%
  \autor{Pedro Nuno Verde Cerqueira}{1260443@isep.ipp.pt}%
  \autor{Rúben Alexandre Abreu Freitas}{1231267@isep.ipp.pt}%
  \autor{Rúben Carlos Zamith Ribeiro}{1161532@isep.ipp.pt}%
}

% Nomes usados nos metadados do PDF (propriedades do ficheiro).
\newcommand{\autoresmeta}{%
  Ana Salomé Ferreira Teixeira;
  Diogo Miguel Gonçalves Ribeiro;
  João Gabriel Domingos de Azevedo;
  Pedro Nuno Verde Cerqueira;
  Rúben Alexandre Abreu Freitas;
  Rúben Carlos Zamith Ribeiro%
}

% ---------------------------------------------------------------
% Docência
%
% Mesma mecânica da autoria: uma linha \docente{nome}{email} por
% elemento.
% ---------------------------------------------------------------
\providecommand{\docente}[2]{}

\newcommand{\listadocentes}{%
  \docente{TODO: nome do docente de ENGCIA}{}%
}

\newcommand{\datarelatorio}{TODO: data de entrega}
